\documentclass[main]{subfiles}
\usepackage[sols]{workbook}
\numbering{ws3-11}

\begin{document}

\ifSubfilesClassLoaded{%
  \chaptermark{\chapterthree}
}{}

\section{Shapes of Solution Trajectories} \label{ws3-11}

\begin{framed}
\centerline{\textbf{Lesson Goals}}

You should be able to:
\begin{itemize}
\item Determine the shape of solution trajectories for a given system of first-order differential equations, and use this information to create a phase portrait for the system
\item Use phase plane analysis and the phase portrait to sketch the solution functions for a system of first-order differential equations
\item Recognize what information we can/cannot obtain from the phase plane, the phase portrait, and the actual solution curves
\end{itemize}
\end{framed}

\begin{enumerate}

  \item 
  \begin{problem}[1in]
  Consider the system:
  \[\left\{\begin{aligned}
  \frac{dx}{dt} &= -x - 2y\\
  \frac{dy}{dt} &= 2x - y
  \end{aligned}\right.\]
  If a solution trajectory in the
  $xy$-phase plane passes through the point $(x,y) = (-1,-3)$, what 
  is the slope of the trajectory at that point?
  \textbf{Hint:} How does the answer relate to the values of $\displaystyle \frac{dx}{dt}$ 
  and $\displaystyle \frac{dy}{dt}$ at $(-1,-3)$?
  \end{problem}

  \begin{solution}
  The question asks us for information about $dy/dx$ at the point
  $(x,y) = (-1,-3)$ in the phase plane.  By the chain rule, we know
  that 
  $$
    \frac{dy}{dt} = \frac{dy}{dx}\cdot \frac{dx}{dt},
  $$
  provided that these derivatives exist.  It follows that
  $$
    \frac{dy}{dx} = \frac{dy/dt}{dx/dt},
  $$
  an expression that makes sense except at points that lie on
  $x$-nullclines.  (If a trajectory crosses an $x$-nullcline,
  $dy/dx$ is undefined at the crossing point---the tangent line
  is vertical.)

  In this example, $\frac{dx}{dt} = 7$ and $\frac{dy}{dt} = 1$ at the
  point $(-1, -3)$.  It follows that the slope of a solution trajectory
  that passes through $(-1,-3)$ is $\boxed{\frac{1}{7}}$.
  \end{solution}



%\probonly{\newpage}


  \item 
  In this problem, we'll look at the system 
 \[\left\{\begin{aligned}
  \frac{dx}{dt} &=y\\
  \frac{dy}{dt} &= -4x
  \end{aligned}
\right.\]
  \begin{enumerate}
    \item
    \begin{problem}[\vfill]
    By writing and solving a differential equation for $y$ in terms of
    $x$, determine the exact shape of the solution trajectories (in the
    phase plane) of this system.
    \end{problem}

    \begin{solution}
    We can use the Chain Rule to write a differential equation relating
    $x$ and $y$.  The Chain Rule tells us that $\displaystyle
    \frac{dy}{dx} = \frac{dy/dt}{dx/dt} =
    -\frac{4x}{y}$.\footnote{Technically, we should only say that
    $\frac{dy}{dx} = \frac{dy/dt}{dx/dt}$ when the denominator,
    $\frac{dx}{dt}$, is not zero; otherwise, we will have divided by
    0.  In practice, we don't worry about keeping track of this: when
    $\frac{dx}{dt} = 0$, we are on a nullcline, and it's easier to get
    information about nullclines by looking at the phase plane
    analysis.}  This is a separable differential equation, so we can
    solve by separating variables:
    \begin{align*}
    \int y~dy & = \int - 4x~dx \\
    \frac{1}{2} y^2 & = - 2x^2 + C \\
    y^2 & = -4x^2 + 2C \\
    4x^2 + y^2 & = 2C
    \end{align*}
    (Of course, $2C$ is just an arbitrary constant, so you could give it
    a new name like $A$.)  This tells us that the solution curves are
    ellipses centered at the origin; each ellipse $4x^2 + y^2 = 2C$ is
    twice as high as it is wide. 
    \end{solution}
\pspace{\newpage}
    \item

    \note{Make sure students see how they need both the phase plane
    analysis and the shape of the trajectories to draw the phase
    portrait --- one or the other alone isn't enough information.}
    
    \begin{problem}[\vfill]
    The following diagram shows the sketch you get by doing a phase
    plane analysis of this system.  Using this together with your answer
    to {{\#2}(a)}, sketch a \textit{phase portrait} of the
    system.  This means: sketch several representative solution
    trajectories in the phase plane.

    \includegraphics[width=2in]{ws/images/system-shape-1-analysis}
    \end{problem}

    \begin{solution}
    From {{\#2}(a)}, we know that the trajectories
    are shaped like ellipses centered at the origin, twice as high as
    they are wide.  We know from the phase plane analysis that the
    trajectories travel clockwise around the origin, with the only
    equilibrium point being the origin.  So, here is a phase portrait:
    \begin{center}
    \includegraphics[width=2.5in]{ws/images/system-shape-1-portrait}
    \end{center}      
    \end{solution}
    
  \end{enumerate}
  
  
  %\probonly{\newpage}
  
  \item 
  In this problem, we'll look at the system 
 \[\left\{\begin{aligned}
  \frac{dx}{dt} &=y\\
  \frac{dy}{dt} &= -3y
  \end{aligned}\right.\]
  \begin{enumerate}
    \item
    \begin{problem}[0.75in]
    Find the exact shape of the solution trajectories.
    \end{problem}

    \begin{solution}
    By the Chain Rule, $\displaystyle \frac{dy}{dx} = \frac{-3y}{y} =
    -3$.  This differential equation is particularly easy to solve,
    since the right side does not depend on $y$ at all.  So, we can just
    integrate both sides with respect to $x$ to get $y = \displaystyle
    \int -3~dx = -3x + C$.  Thus, the trajectories are lines of slope
    $-3$.
    \end{solution}

    \item
    \begin{problem}[3in]
    Do a qualitative phase plane analysis of this system.
    \end{problem}

    \begin{solution}
    We'll follow the steps outlined in the previous lesson:
    \begin{enumerate}[start=0,label=\arabic*.]
    \item The equations are already factored.
    \item The only $\frac{dx}{dt} = 0$ nullcline is $y = 0$.
    \item The only $\frac{dy}{dt} = 0$ nullcline is $y = 0$.
    \item So, the equilibrium points are all points with $y = 0$.  We
    have the following picture so far:
    \begin{center}
    \includegraphics[width=2in]{ws/images/system-shape-2-analysis-step3} 
    \end{center}
    \item There are no $\frac{dx}{dt} = 0$ nullclines left to orient.
    \item There are no $\frac{dy}{dt} = 0$ nullclines left to orient.
    \item There are two regions in which we should draw arrows showing
    the general direction of trajectories --- the region above $y = 0$
    and the region below $y = 0$.  For the top region, we can test a
    point like $(1, 1)$, and we find that $\frac{dx}{dt} > 0$ and
    $\frac{dy}{dt} < 0$, so the trajectories go to the right and
    down.  For the bottom region, we can test a point like $(1, -1)$,
    and we find that $\frac{dx}{dt} < 0$ and $\frac{dy}{dt} > 0$, so
    the trajectories go to the left and up.
  \end{enumerate}
  This gives us the following sketch:
  \begin{center}
  \includegraphics[width=2in]{ws/images/system-shape-2-analysis}
  \end{center}
  \end{solution}

  \item
    \note{Students will often draw full lines with slope $-3$, forgetting that the lines should terminate on the $x$-axis based on the phase plane analysis.}
  \begin{problem}[2in]
  Using your answers to {{\#3}(a)} and
  {{\#3}(b)}, draw a phase portrait of the system.
  \end{problem}

  \begin{solution}
  We know that the trajectories are lines of slope $-3$, and we can
  figure out the appropriate directions from our phase plane
  analysis.  So, the phase portrait looks like this:
  \begin{center}
  \includegraphics[width=2in]{ws/images/system-shape-2-portrait}
  \end{center}
  (One trajectory is shown in red for the next few parts.)
  \end{solution}

  \item
  \begin{problem}[0.75in]
  Let's focus on one particular solution, the one satisfying the
  initial conditions $x(0) = 0$ and $y(0) = 6$.  What is the exact
  shape of the solution trajectory with these initial conditions?
  Write an equation which describes the shape.
  \end{problem}

  \begin{solution}
  From {{\#3}(a)}, we know that all of the
  trajectories have shapes given by $y = -3x + C$.  We are interested
  in the particular trajectory containing the point $(0, 6)$, so we
  simply plug in $x = 0$, $y = 6$ to solve for $C$.  This gives $C =
  6$, so the shape of the particular trajectory we care about is
  $\boxed{y = -3x + 6}$.  (This is the trajectory highlighted in red
  in the previous picture.)

  Notice that the trajectory goes along the line $y = -3x + 6$, but it
  does not go along the entire line, only the part above $y = 0$.
  This is something we only know because of the phase plane analysis. 
  \end{solution}


  \item
  \note{Students often think that $x(t)$ should progress linearly from 0 to 2, terminating at some unknown time $t$, and similarly for $y(t)$ progressing from 6 to 0. I like to let them make and correct this mistake themselves by continuing to work through parts (f) and (g). This is also a good time to remind them that when we draw our phase planes and portraits, we only know the general direction (and perhaps the shape) of the trajectories, but not how much time it takes to get anywhere.}
 
  \begin{problem}[0.5in]
  On the following set of axes, sketch the graphs of $x(t)$ and $y(t)$
  if $x(0) = 0$ and $y(0) = 6$.

  \probonly{\includegraphics[height=2in]{ws/images/romeo-juliet-time-axes}}
  \end{problem}

  \begin{solution}
  If we start at $(0, 6)$ on the red trajectory, then we go down and
  to the right (so $x$ increases and $y$ decreases) toward $(2, 0)$.
  So, here are reasonable graphs of $x(t)$ and $y(t)$:
  \begin{center}
  \includegraphics[width=2in]{ws/images/system-shape-2-sketch-x-and-y}
  \end{center}
  Note that, if we want to incorporate everything we know so far, we
  should use the fact that $y(t) = -3x(t) + 6$, which we showed in the
  previous part.  However, it's not so easy to portray that
  relationship in our graphs --- that's one reason looking at the
  phase portrait is useful.
  \end{solution}
\pspace{\newpage}
  \item
  \note{It is fairly easy for them to screw up the constants here.  (For
  instance, once they find that the solution is $x(t) = -
  \frac{C_1}{3} e^{-3t} + C_2$ and $y(t) = C_1 e^{-3t}$, they may want
  to rewrite $x(t)$ as $C_1 e^{-3t} + C_2$ without changing $y(t)$.)}

  \begin{problem}[2in]
  In fact, you know enough to solve the system 
 \[\left\{\begin{aligned}
  \frac{dx}{dt} &=y\\
  \frac{dy}{dt} &= -3y
  \end{aligned}\right.\]
  Find the general solution of the system; that is, find the general solution for $x(t)$ and $y(t)$. 
\textbf{Hint:} First solve for $y(t)$.
  \end{problem}

  \begin{solution}
  First, remember what it means to solve the system: we are looking
  for two functions, $x(t)$ and $y(t)$, with the property that $x'(t)
  = y(t)$ and $y'(t) = -3y(t)$.  Let's start with the second equation,
  which looks simpler because it's really just about a single function
  $y(t)$; this equation says that $y(t)$ is a function of $t$ whose
  derivative is $-3$ times the function itself.

  You should be exceptionally familiar with this particular
  differential equation by now; we have seen it many times.  As we
  learned a while ago, the general solution
  of $\frac{dy}{dt} = -3y$ is $y(t) = C_1 e^{-3t}$, where $C_1$ is any
  constant.

  Now, let's go back and look at the other equation, $x'(t) = y(t)$.
  We now know that $y(t) = C_1 e^{-3t}$, so we are looking for a
  function $x(t)$ whose derivative (with respect to $t$) is $C_1
  e^{-3t}$.  That is, we are trying to solve the differential equation
  $x'(t) = C_1 e^{-3t}$.  Since the right side is purely in terms of
  $t$, we can just integrate with respect to $t$ to get $x(t) =
  \displaystyle \int C_1 e^{-3t}~dt = - \frac{C_1}{3} e^{-3t} + C_2$.
  So, we've solved the system: $\boxed{x(t) = - \frac{C_1}{3} e^{-3t}
  + C_2 \textrm{ and } y(t) = C_1 e^{-3t}}$.
  \end{solution}

  \item

  \note{Another good thing to discuss here is that $x(t)$ and $y(t)$
  really do satisfy the equation they found in
  {{\#3}(d)}.}

  \begin{problem}[1in]
  Find the solution of the system satisfying $x(0) = 0$ and $y(0) =
  6$.  Does your answer match the graphs you sketched in
  {{\#3}(e)}? 
  \end{problem}

  \begin{solution}
  We found in the previous part that the general solution to the
  system was $x(t) = - \frac{C_1}{3} e^{-3t} + C_2$ and $y(t) = C_1
  e^{-3t}$.  As always, if we know the general solution and want to
  find a particular solution, all we really need to do is to plug in
  the initial conditions to solve for the constants $C_1$ and $C_2$.

  Here, the initial condition $x(0) = 0$ tells us that $-
  \frac{C_1}{3} + C_2 = 0$, so $C_2 = \frac{C_1}{3}$.  The initial
  condition $y(0) = 6$ tells us that $C_1 = 6$.  So, $C_2 =
  \frac{6}{3} = 2$.  Thus, the solution we're looking for is
  $\boxed{x(t) = 2 - 2e^{-3t}$, $y(t) = 6 e^{-3t}}$.  (We can easily
  verify that this satisfies the original system of differential
  equations, too.)  This does match the graphs we sketched.

  Notice that it also matches the relationship $y(t) = - 3x(t) + 6$
  that we came up with in
  {{\#3}(d)}. 
  \end{solution}

\end{enumerate}


\pspace{\newpage}

\item  \label{SIR}

    Kermack and McKendrick (1927) used systems of differential equations 
    to model the spread of an infectious 
    disease within a population.  They modeled interactions between
    susceptible ($S$) and infected ($I$) individuals in much the same way that
    Lotka and Volterra modeled interactions between predators and prey,
    and proposed the following system:
         \begin{align*}
\frac{dS}{dt} &= -\alpha S I  \\
\frac{dI}{dt} &= \alpha S I - \beta I
  \end{align*}
    where $\alpha$ and $\beta$ are positive constants related to the
    transmission rate and recovery rate, respectively.

  \begin{enumerate}
  \item
  \begin{problem}[0.5in]
    If the infected population were quarantined (effectively changing
    $\alpha$ to zero), what would the model predict?
  \end{problem}

  \begin{solution}
    If the infected population were quarantined, then $dI/dt = -\beta I$.
    This implies that the infection would die off exponentially.  Thinking
    optimistically, let's assume this means that the infected population
    is gradually recovering (as opposed to dying!) from the infection.
    After the infection has run its course, the infected population is
    immune, meaning that they're not returned to the susceptible population.
  \end{solution}

  \item
  \begin{problem}[3in]
    Perform a phase plane analysis for this system, assuming $\alpha = 0.01$
    and $\beta = 0.1$. Let $S$ represent the $x$-axis and $I$ represent the $y$-axis.
  \end{problem}

  \begin{solution}
    Here are the model equations with the given choices of parameters:
        \begin{align*}
\frac{dS}{dt} &= -0.01 S I  \\
\frac{dI}{dt} &= 0.01 S I - 0.1 I=0.01I(S-10)
  \end{align*}
    There are two $S$-nullclines:  $S = 0$ and $I = 0$.  This makes sense 
    given the context:  the susceptible population cannot change if either
    (i) there are no susceptible people or (ii) there is no infection.  There
    are also two $I$-nullclines:  $I = 0$ and $S = 10$.  Interestingly, $I = 0$
    is {\bf both} an $S$-nullcline and an $I$-nullcline.  This means that
    every single point on the $I = 0$ axis is an equilibrium point!  Again, this
    makes perfect sense:  If there is no infection, then the susceptible 
    population (no matter what its value) would remain constant.

    Below is a sketch of a phase plane diagram for this system, where we
    have elected to let $I$ correspond to the vertical axis.  We have drawn
    lots of bold dots along the horizontal axis in order to indicate that
    every point on that axis is an equilibrium.
    \begin{center}
    \includegraphics[width=2.5in]{ws/images/SIR.eps}
    \end{center}

  \end{solution}


  \item
  \begin{problem}
    Again using $\alpha = 0.01$ and $\beta = 0.1$, determine the exact 
    shape of the solution trajectories.  Then (using a graphing calculator
    if necessary), draw the trajectory corresponding to the initial conditions
    $S(0) = 20$, $I(0) = 5$.  Interpret your result in an epidemiological
    context.
  \end{problem}

  \begin{solution}
    We know that $dI/dS$ is undefined along the $S$-nullclines, but elsewhere
    $$
	\frac{dI}{dS} = \frac{dI/dt}{dS/dt} \; = \; \frac{0.01 I(S-10)}{-0.01 S I} \; = \; \frac{10 - S}{S} \; = \; \frac{10}{S} - 1.
    $$
    The solutions of this differential equation tell us the shapes of the
    curves that trajectories follow in the $SI$-plane.  To solve this
    differential equations, we merely need to integrate with respect to $S$:
    $$
      I = 10 \ln(S) - S + C.
    $$
    Note:  Absolute value bars are unnecessary in the $\ln(S)$ term, since $S$
    represents a population and therefore cannot be negative.  For the given
    initial conditions, we can compute the value of the constant of integration:
    $$
      5 = 10 \ln(20) - 20 + C \qquad \implies \qquad C = 25 - 10 \ln(20).
    $$
    It follows that the solution trajectory satisfying the given initial conditions
    lies along the curve
    $$
      I = 10 \ln(S) - S + 25 - 10 \ln(20)
    $$
    in the phase plane.  The motion along the curve must be in the direction of
    decreasing $S$, bearing in mind that $dS/dt$ is never positive.  A sketch of
    the solution trajectory appears in green in the following figure.

    \begin{center}
    \includegraphics[width=2.5in]{ws/images/SIR-traj.eps}
    \end{center}


    The model predicts that the epidemic will get worse before it gets better:
    the number of infected individuals will increase for awhile.  Once $S$
    decreases to $10$, the contact rate between susceptible and infected 
    individuals is diminished, due to lower availability of susceptible people.
    Then, the infected population begins to diminish.  The model also yields
    some good quantitative predictions:  The infected population will peak when
    $S = 10$, at which point $I \approx 8.069$.  Moreover, if we use a computer
    to estimate where the trajectory crosses the $S$-axis, we find that when
    the infection has died off, there is a remaining susceptible population of
    approximately $2.006$.
  \end{solution}

  \end{enumerate}


\end{enumerate}



\end{document}